\documentclass[11pt]{article}
\usepackage[a4paper,margin=22mm]{geometry}
\usepackage{amsmath,amssymb,mathtools,bm}
\usepackage{booktabs,array,microtype,xurl,hyperref}
\hypersetup{colorlinks=true,linkcolor=blue,citecolor=blue,urlcolor=blue,
 pdftitle={P54PQ-v2 bosonic Coleman-Weinberg matching Hessian}}
\setlength{\emergencystretch}{1em}
\newcommand{\tr}{\operatorname{tr}}
\newcommand{\diag}{\operatorname{diag}}
\newcommand{\MSbar}{\overline{\mathrm{MS}}}
\newcommand{\PiB}{\Pi_{B}^{\rm hard}}
\title{P2 Bosonic Coleman--Weinberg Matching Hessian\\
for the Hierarchical \(P54PQ\)-v2 Vacuum}
\author{Route-F four-dimensional mainline}
\date{2 September 2026}

\begin{document}
\maketitle

\begin{abstract}
We compute the bosonic one-loop zero-momentum matching curvature on the real
four-plane of the tuned Standard-Model Higgs doublet.  The declared scheme is
background-field Landau gauge and \(\MSbar\) at
\(\mu_U=g_U\omega\).  A hard/soft spectral split retains all 290 massive
physical scalar directions and 33 massive vector directions while leaving
the 38 scalar and 12 vector infrared degrees of freedom in the lower-energy
theory.  A matrix-function Fr\'echet formula differentiates the complete
field-dependent \(328\times328\) scalar Hessian without subtracting nearly
equal one-loop potentials.  At the fixed point
\(\sigma/\omega=0.1265\), the result is
\[
 \PiB/\omega^2=-0.0202366510950\,\bm 1_4,
 \qquad \delta\xi_{02,B}=-0.0202367022071.
\]
This is a scheme-dependent matching coefficient, not a gauge-independent
pole mass.  The unfit heavy-fermion projection is retained explicitly as
\(\eta_Y(\mu)\), so the renormalization condition is
\(\delta\xi_{02}=[\kappa_B+\eta_Y]/w_{10}\).
\end{abstract}

\tableofcontents

\section{Background and the precise claim}

The input is the previously verified hierarchical stationary point
\begin{equation}
 \frac{\sigma}{\omega}=0.1265,\qquad
 \frac{v_s}{\omega}=0.25,\qquad
 g_U=0.562602889985,
 \label{eq:background}
\end{equation}
with tree parameter \(\xi_{02}=0.147371094023\).  Its complete real scalar
Hessian has dimension
\[
 54+2(126)+2(10)+2=328.
\]
There are 38 zero directions: 33 gauge-orbit directions, one physical PQ
Goldstone, and a real four-plane \(L\) carrying one
\((\bm1,\bm2,|Y|=1/2)\) multiplet.  There are no negative modes.  The next
doublet has
\begin{equation}
 \Delta_D=0.0305570302561\,\omega^2.
 \label{eq:gap}
\end{equation}

The calculation below determines the heavy-boson part of the one-loop
zero-momentum matching matrix on \(L\).  It does \emph{not} determine a
gauge-independent pole mass.  The latter additionally requires the
momentum-dependent self-energy, wave-function matching, a consistent
counterterm prescription, and the fitted heavy-Yukawa matrices.  This
distinction is compulsory because an effective potential away from a
stationary point is gauge dependent; the Nielsen identity does not turn an
arbitrary Hessian entry into an observable~\cite{Nielsen1975}.

\section{Gauge, ghost, and infrared prescription}

We use background-field \(R_\xi\) gauge at
\begin{equation}
 \xi_{\rm gf}=0
 \qquad\text{(Landau gauge)}
 \label{eq:landau}
\end{equation}
and \(\MSbar\).  At constant background the hard bosonic functional is
\begin{align}
 V_{1,B}^{\rm hard}(x;\mu)
 =\frac{1}{64\pi^2}\bigg[&
 \tr_{290} f_s\!\left(M_s^2(x)\right)
 +3\tr_{33} f_v\!\left(M_v^2(x)\right)\bigg],
 \label{eq:cw}\\
 f_s(z)&=z^2\left(\log\frac{z}{\mu^2}-\frac32\right),
 &f_v(z)&=z^2\left(\log\frac{z}{\mu^2}-\frac56\right).
 \label{eq:fdefs}
\end{align}
The vector coefficient three and the \(5/6\) finite constant are the usual
Landau-\(\MSbar\) vector result~\cite{ColemanWeinberg1973,Martin2001}.  The
Faddeev--Popov eigenvalues are
\(m_{\rm gh}^2=\xi_{\rm gf}M_v^2=0\); hence there is no hard ghost term in
this declared scheme.

The subscripts on the traces are a Wilsonian spectral split.  The 290
positive scalar eigenvalues are retained; the 38 zero modes are assigned to
the infrared theory.  Similarly, the 33 positive vector eigenvalues are
retained and the 12 unbroken vectors are infrared fields.  This split is
locally unambiguous because
\begin{equation}
 \lambda_{s,\min}^{\rm hard}=0.0135204700328\,\omega^2,
 \qquad
 \lambda_{v,\min}^{\rm hard}=0.01600225\,g_U^2\omega^2,
 \label{eq:hardgaps}
\end{equation}
whereas all soft eigenvalues vanish at the matching background.  Hard--soft
mixing is retained in the divided-difference term below; ``remove the soft
trace'' therefore does not mean projecting every derivative onto the hard
subspace.

\section{First-principles second variation}

Let \(z^a\), \(a=1,\ldots,4\), be orthonormal canonical coordinates on
\(L\).  For either bosonic mass operator, write
\begin{equation}
 M^2(z)=M_0^2+z^a A_a+\frac12z^az^b B_{ab}+O(|z|^3).
 \label{eq:massjet}
\end{equation}
If \(M_0^2=U\diag(\lambda_i)U^T\), define all matrices below in this
eigenbasis.  Differentiating \(\tr f(M^2)\) once gives
\begin{equation}
 \partial_a\tr f(M^2)=\tr\!\left[f'(M^2)\partial_aM^2\right].
 \label{eq:first}
\end{equation}
The Fr\'echet derivative of \(f'\) has components
\begin{equation}
 \left[D f'(M_0^2)[A_b]\right]_{ij}
 =f'^{[1]}(\lambda_i,\lambda_j)(A_b)_{ij},
 \label{eq:frechet}
\end{equation}
where
\begin{equation}
 f'^{[1]}(x,y)=
 \begin{cases}
 \dfrac{f'(x)-f'(y)}{x-y},&x\ne y,\\[6pt]
 f''(x),&x=y.
 \end{cases}
 \label{eq:divdiff}
\end{equation}
Therefore the exact second variation is
\begin{equation}
 \boxed{
 \partial_a\partial_b\tr f(M^2)=
 \sum_i f'(\lambda_i)(B_{ab})_{ii}
 +\sum_{i,j}f'^{[1]}(\lambda_i,\lambda_j)
 (A_a)_{ij}(A_b)_{ji}.}
 \label{eq:master}
\end{equation}
For a soft eigenvalue we set \(f'=f''=0\), the derivative of the locally
flat infrared subtraction.  Equation~\eqref{eq:master} still includes the
hard--soft terms because their divided differences are nonzero.

For \(f_c(z)=z^2[\log(z/\mu^2)-c]\), the derivatives used in the code are
\begin{align}
 f_c'(z)&=z\left[2\log\frac{z}{\mu^2}-2c+1\right],\nonumber\\
 f_c''(z)&=2\log\frac{z}{\mu^2}-2c+3.
 \label{eq:fderivatives}
\end{align}
The scalar mass operator is not reconstructed from a tabulated spectrum:
it is the automatic second derivative of the full tensor potential,
\begin{equation}
 M_s^2(x)=\frac{\partial^2V_{\rm tree}}
 {\partial x^I\partial x^J}.
 \label{eq:scalaroperator}
\end{equation}
Because the action is quartic, \(M_s^2(x)\) is exactly quadratic in \(x\).
Thus centered evaluations at \(\epsilon/\omega=0.02\) determine \(A_a\)
and \(B_{aa}\), and a combined displacement determines \(B_{01}\).
Repeating one direction at \(\epsilon/\omega=0.01\) gives relative residuals
\begin{equation}
 r_A=2.32\times10^{-14},\qquad r_B=2.15\times10^{-11}.
 \label{eq:steps}
\end{equation}

For vectors the derivation is analytic.  If
\(O_A(x)=T_Ax\) is the gauge-orbit matrix, then
\begin{equation}
 M_v^2(x)=g_U^2O(x)^TO(x),
 \label{eq:vectormass}
\end{equation}
so, with \(O_a=Tq_a\),
\begin{align}
 A_a^{(v)}&=g_U^2(O_a^TO_0+O_0^TO_a),\nonumber\\
 B_{ab}^{(v)}&=g_U^2(O_a^TO_b+O_b^TO_a).
 \label{eq:vectorjet}
\end{align}

\section{Light-doublet basis and symmetry reduction}

Let \(Z\) be the tree zero-eigenspace and let \(S\) contain the 33 gauge
directions plus the physical PQ direction.  The numerical light projector is
obtained from
\begin{equation}
 L=\operatorname{range}\!\left[(1-SS^+)Z\right].
 \label{eq:lightprojector}
\end{equation}
Its rank is four and its Casimirs are
\(C_3=0\), \(C_2=3/4\), and \(Y^2=1/4\).  The hard spectral traces commute
with the unbroken Standard Model.  The realification of this charged complex
doublet has a unique invariant symmetric bilinear form.  Consequently
\begin{equation}
 \PiB\big|_L=\kappa_B\bm1_4.
 \label{eq:schur}
\end{equation}
This is a symmetry statement, not a numerical fit.  Numerically, two
independent diagonal directions agree to \(2.29\times10^{-13}\), the tested
mixed component is \(7.70\times10^{-14}\omega^2\), and the complete vector
block has isotropy residual \(2.94\times10^{-18}\).

\section{Numerical result and scale dependence}

At \(\mu_U=g_U\omega\), Eq.~\eqref{eq:master} gives
\begin{align}
 \frac{\Pi_s^{\rm hard}}{\omega^2}
 &=-0.0182525705951\,\bm1_4,\nonumber\\
 \frac{\Pi_v^{\rm hard}}{\omega^2}
 &=-0.00198408049986\,\bm1_4,\nonumber\\
 &\boxed{\frac{\PiB}{\omega^2}
 =-0.0202366510950\,\bm1_4}
 \label{eq:result}
\end{align}
The off-diagonal entries of the directly assembled vector matrix are below
\(1.1\times10^{-18}\omega^2\).  The four equal components in
Eq.~\eqref{eq:result} are fixed by Eq.~\eqref{eq:schur}; only two diagonal
directions and one mixed direction need expensive scalar evaluations.

The coefficient is scale dependent, as a matching coefficient must be:
\begin{center}
\begin{tabular}{@{}rrrrr@{}}
\toprule
\(\mu/\mu_U\)&\(\kappa_s/\omega^2\)&\(\kappa_v/\omega^2\)&
\(\kappa_B/\omega^2\)&\(\delta\xi_{02,B}\)\\
\midrule
0.5& \( 0.0107814599\)&\( 0.0060371453\)&\( 0.0168186052\)&\( 0.0168186477\)\\
1.0& \(-0.0182525706\)&\(-0.0019840805\)&\(-0.0202366511\)&\(-0.0202367022\)\\
2.0& \(-0.0472866011\)&\(-0.0100053063\)&\(-0.0572919074\)&\(-0.0572920521\)\\
\bottomrule
\end{tabular}
\end{center}
The large finite-scale movement is not an uncertainty band.  It is cancelled
order by order by the running renormalized parameter and by the scale
dependence of the fermionic matching contribution.  Quoting only the middle
number as a physical prediction would therefore be incorrect.

\section{Retuning and the heavy-Yukawa nuisance}

The normalized light mode has
\begin{equation}
 w_{10}=\langle h|P_{10}|h\rangle=0.999997474290.
 \label{eq:w10}
\end{equation}
Since the tree Hessian responds to \(\delta\xi_{02}\) as
\(-\delta\xi_{02}P_{10}\), the bosonic counterterm required to keep the
doublet massless at the matching surface is
\begin{equation}
 \delta\xi_{02,B}(\mu_U)
 =\frac{\kappa_B(\mu_U)/\omega^2}{w_{10}}
 =-0.0202367022071.
 \label{eq:bosretune}
\end{equation}

The heavy-fermion contribution is deliberately not guessed.  Define the
matching nuisance
\begin{equation}
 \eta_Y(\mu)=\frac{1}{\omega^2}
 \langle h|\Pi_{\rm heavy-Y}(0;\mu)|h\rangle.
 \label{eq:etay}
\end{equation}
The complete matching relation carried into the global flavor analysis is
then
\begin{equation}
 \boxed{\delta\xi_{02}(\mu)=
 \frac{\kappa_B(\mu)/\omega^2+\eta_Y(\mu)}{w_{10}}.}
 \label{eq:fullretune}
\end{equation}
No value or prior for \(\eta_Y\) is fabricated here.  The independent
one-light-doublet safety gate remains
\begin{equation}
 \left\|Q\left[\Pi_B+\Pi_Y-\delta\xi_{02}P_{10}\right]Q\right\|
 <\Delta_D,
 \qquad Q=1-|h\rangle\langle h|,
 \label{eq:safety}
\end{equation}
with \(\Delta_D\) from Eq.~\eqref{eq:gap}.  A global heavy-Yukawa fit must
evaluate Eq.~\eqref{eq:safety}; the present calculation closes the bosonic
matching coefficient, not that later flavor-dependent norm bound.

\section{Reproducibility and promotion boundary}

The machine calculation is
\begin{verbatim}
route_f/code/verify_p54_p2_bosonic_cw.py
route_f/output/p54_p2_bosonic_cw.json
route_f/output/p54_p2_bosonic_cw.md
\end{verbatim}
It evaluates eight complete field-dependent scalar Hessians and passes 14 of
14 physics/numerical checks.  The soft--soft first-derivative norms are
\(3.88\times10^{-15}\) and \(3.08\times10^{-15}\), closing the declared IR
split.  An independent random-matrix regression of Eq.~\eqref{eq:master}
against a direct four-point finite difference has relative residual
\(2.47\times10^{-9}\).

The P2 bosonic matching task is therefore closed in the declared scheme.
The full pole mass remains unpredicted, and \(\eta_Y\) is now an explicit
matching nuisance rather than a hidden assumption.  This is the appropriate
boundary at which P3 may begin a global flavor/seesaw fit: P3 must fit or
profile \(\eta_Y\), propagate it through Eq.~\eqref{eq:fullretune}, and test
Eq.~\eqref{eq:safety}.

\begin{thebibliography}{9}
\bibitem{ColemanWeinberg1973}
S.~Coleman and E.~Weinberg,
``Radiative Corrections as the Origin of Spontaneous Symmetry Breaking,''
Phys. Rev. D \textbf{7} (1973) 1888.

\bibitem{Nielsen1975}
N.~K. Nielsen,
``On the Gauge Dependence of Spontaneous Symmetry Breaking in Gauge
Theories,'' Nucl. Phys. B \textbf{101} (1975) 173.

\bibitem{Martin2001}
S.~P. Martin,
``Two-loop effective potential for a general renormalizable theory and
softly broken supersymmetry,'' Phys. Rev. D \textbf{65} (2002) 116003,
arXiv:hep-ph/0111209.
\end{thebibliography}

\end{document}
